\section{多次元尺度構成法}
\subsection{授業内容}
\subsubsection{科目の中でのこのコマの位置づけ}
数量化まで学ぶことによって，カテゴリに適切な数値を割り振るという尺度化の原点に立ち返ることができた。
ここではさらに，心理尺度のような回答法ではない変数間同士の関係から，次元を取り出して分析する多次元尺度構成法について考える。
この方法は実験刺激や知覚的反応，直感的判断などを対象にできるため，応用可能性が非常に高いだろう。

多次元尺度構成法を理解するためには，まず実際の距離行列を分析して地図を再構成できるかどうかを見るところから始めるのが良いだろう。データとして与えられるのが距離行列であり，行動計量学ではこれを心理的な距離や意味的な距離が数値化されたものだと捉えることで，心理学的な地図を作っていると解釈してきた。この仮定には最大限の注意を払いつつ，必ずしも計量的でない場合の数値化をする非計量的多次元尺度法に拡張することで，更なる心理学的用途が広がることを見る。なお，多次元尺度方は数量化IV類と同じである。

多次元尺度法は分析の元が距離行列であり，その基底を固有値分解によって得るといういみで，多変量解析としてはお馴染みの考え方であるともいえる。しかしデータが距離(を意味するもの)であれば良く，数学的にも簡単な拡張をすることで，個人差を表現するモデルに拡張することもできる。また心理尺度に対する考え方として，個人の内的な次元からの近さに応じて反応すると考える，展開型のモデルを使うことは，心理尺度の利用に新な視点をもたらす。

\subsubsection{コマ主題細目}
\begin{description}
    \item[多次元尺度構成法] 多次元尺度構成法(MDS)は，距離行列を元にした多変量解析の一種であり，距離関係から次元(基底)を選び出し，対象に座標を与える方法である。まずは座標の復元例から考え，抽出する次元数をどのようにして求めるかといった基礎的な知識を得る。また，行列の考え方からみると正方対称行列の固有値分解であるから，これを確認するだけでも他の多変量解析と合わせた統合的な理解ができるだろう。因子分析モデルも多次元尺度法の一種であるということもできる。



    \item[距離と心理学のデータ] 距離行列があれば次元が抽出できることが分かったが，さて何を距離とみなすかを考えれば，非常に多くの可能性が広がることがわかる。距離の定義は非負で対称性と三角不等式が成り立つことであり，様々な距離の定め方があるし，共分散や相関もその一種と考えることができる。元になるデータも尺度評定を用いる方法，刺激の混同率，代替価/連想価，刺激の汎化勾配，    反応潜時，ソシオメトリックなデータなど，心理学の様々な領域で得られるデータが，距離とみなすことができる。応用可能な領域が広いことを知ることで，統計モデルをハンドリングできることになる。
        \begin{flushright}
            $\to$データの例に関しては\citet{高根芳雄1980-02}のPp.14--27.
        \end{flushright}
    \item[非計量多次元尺度法] 計量MDSによって算出される座標は元のデータをうまく復元するが，心理学的なデータの場合はデータの大小関係の表現(順序尺度水準)がせいぜいであり，これに対応した非計量多次元尺度法が考案されている。この手法を用いることで，一対比較や順序比較などのより制限の少ないデータからであっても数量関係を導き出すことができる。
        \begin{flushright}
            $\to$ 計量・非計量多次元尺度構成方については，すでに絶版になったが\citet{岡太彬訓1994-01-15}がもっとも簡潔でわかりやすく説明している。手に入るところでは\citet{足立浩平2006-07}のP.135--143，あるいは\citet{kosugi2019b}のP.199-203.
        \end{flushright}
        \item[多次元尺度構成法の展開]多次元尺度構成法で作られたものは地図である。地図には点を書き込んだり，複数の地図を重ねたりできるように，多次元尺度法で得られた地図にも情報を追加したり，モデルを展開するなどして様々な応用的モデルを作ることができる。ここではPrefmapや楕円モデル(非対称MDS)，INDSCALなど応用例をいくつか示し，この技術の応用可能性を考える。
    \item[展開法] Coombsが考えた心理尺度の展開法は，サーストン法やリッカート法とはまた別の尺度化の考え方を表している。この方法は被験者とカテゴリーの両方を地図上にプロットできる。具体的な分析例をみながら，尺度やそれに数字を与える方法についての考え方を見る。
        \begin{flushright}
            $\to$展開型モデルについては，多少複雑な工夫が組み込まれているが，\citet{shimizu2018}が参考になる。
        \end{flushright}


\end{description}
\subsubsection{キーワード}
\begin{itemize}
    \item 多次元尺度構成法
    \item 非計量多次元尺度構成法
    \item 個人差多次元尺度構成法
    \item 展開型多次元尺度法
\end{itemize}
\subsection{授業情報}
\paragraph{コマの展開方法}
講義
\subsubsection{予習・復習課題}
\paragraph{予習}数量化の関係を再確認しよう。すなわち数値に値を与えるというものであり，尺度のカテゴリだけでなくより一般的な心理的刺激を考え，どういったモデルで表現できるかを考えると本講だけでなく後期の授業にもつながる気づきを得るだろう。
\paragraph{復習}多次元尺度法によって，どのような分析ができるかを考えてみよう。とくに尺度法にかかわらず，実験刺激からの反応を距離と見做せる関係にすれば分析でき，その結果をどのように解釈するかについても自由度はかなり多い。紹介された発展的なモデルなどについても自分の研究関心にどのように応用できるか考えてみよう。